% ======================================================================
%  PROPOSTA PRELIMINAR DE TCC
%  Engenharia de Controle e Automação - IFSP / Campus São José dos Campos
%  Aluno: Gustavo Henrique Ferraioli
%  Tema: Controlador de atitude para quadrirrotor embarcado no ArduPilot,
%        validado em simulação SITL (Software-In-The-Loop)
%  Modelo baseado na "Proposta de TCC - Guilherme - Revisado"
% ======================================================================
\documentclass[12pt,a4paper]{article}

% ---------- Pacotes ----------
\usepackage[utf8]{inputenc}
\usepackage[T1]{fontenc}
\usepackage[brazil]{babel}
\usepackage[a4paper,left=3cm,right=2cm,top=3cm,bottom=2cm]{geometry}
\usepackage{indentfirst}      % indenta o primeiro parágrafo de cada seção
\usepackage{setspace}         % espaçamento entre linhas
\usepackage{enumitem}         % listas customizadas
\usepackage{titlesec}         % formatação dos títulos de seção
\usepackage{graphicx}
\usepackage{array}            % colunas customizadas nas tabelas
\usepackage{amsmath}          % equações
\usepackage{hyperref}
\hypersetup{hidelinks}

% ---------- Espaçamentos ----------
\onehalfspacing               % 1,5 linha
\setlength{\parindent}{1.25cm}
\setlength{\parskip}{0.3em}

% ---------- Estilo das seções (numeradas, em negrito) ----------
\titleformat{\section}{\normalfont\bfseries}{\thesection.}{0.5em}{}
\titleformat{\subsection}{\normalfont\bfseries}{\thesubsection.}{0.5em}{}
\titlespacing*{\section}{0pt}{1.2em}{0.6em}
\titlespacing*{\subsection}{0pt}{0.8em}{0.4em}

% ---------- Tipo de coluna para tabelas ----------
\newcolumntype{L}[1]{>{\raggedright\arraybackslash}p{#1}}

% ======================================================================
\begin{document}

% ----------------------------------------------------------------------
%  CABEÇALHO
% ----------------------------------------------------------------------
\begin{center}
    {\large\bfseries ENGENHARIA DE CONTROLE E AUTOMAÇÃO}\\[0.3em]
    (Trabalho de Conclusão de Curso -- 2026)\\[1.2em]
    {\large\bfseries PROPOSTA PRELIMINAR DE TCC}\\[1.8em]
\end{center}

\noindent\textbf{Aluno:} Gustavo Henrique Ferraioli\\[0.4em]
\noindent\textbf{Orientador:} Prof. Clécio Fischer\\[0.2em]
\noindent\textbf{Coorientador:} [A definir]\\[1.5em]

\begin{center}
    {\large\bfseries PROJETO E IMPLEMENTAÇÃO EMBARCADA DE UM\\
    CONTROLADOR DE ATITUDE PARA QUADRIRROTOR NO\\
    FIRMWARE ARDUPILOT, VALIDADO EM SIMULAÇÃO\\
    \emph{SOFTWARE-IN-THE-LOOP}}
\end{center}
\vspace{1.5em}

% ----------------------------------------------------------------------
%  1. ÁREA DE INTERESSE
% ----------------------------------------------------------------------
\section{Área de Interesse}

\begin{itemize}[leftmargin=2cm,nosep,topsep=2pt]
    \item Modelagem e Simulação de Sistemas Dinâmicos Não-Lineares
    \item Sistemas de Controle Moderno (realimentação de estados e controle ótimo)
    \item Controle Digital e Implementação Embarcada de Leis de Controle
    \item Identificação de Sistemas
    \item Modelagem de Distúrbios Atmosféricos e Rejeição de Perturbações
    \item Sistemas Embarcados de Tempo Real Aplicados a Veículos Aéreos
          Não Tripulados (VANTs)
\end{itemize}

% ----------------------------------------------------------------------
%  2. MOTIVAÇÃO
% ----------------------------------------------------------------------
\section{Motivação}

Os veículos aéreos não tripulados de asas rotativas, em especial os
quadrirrotores, deixaram de ser objeto exclusivo de pesquisa e passaram a ocupar
aplicações cotidianas de inspeção de estruturas, agricultura de precisão,
levantamento topográfico, monitoramento ambiental e logística de última milha. Essa
difusão foi possibilitada tanto pela redução de custo de sensores inerciais quanto
pela maturidade dos pilotos automáticos de código aberto, que hoje entregam, em
plataformas de baixo custo, funcionalidades antes restritas a sistemas
aeronáuticos certificados.

Do ponto de vista da Engenharia de Controle e Automação, essa popularização é
notável por um motivo específico: o quadrirrotor é uma planta \emph{subatuada} ---
possui seis graus de liberdade e apenas quatro entradas independentes ---,
\emph{não-linear}, \emph{fortemente acoplada} entre os eixos de atitude e de
translação, e \emph{instável em malha aberta}. Sem realimentação, a aeronave
simplesmente não voa. Todo o comportamento observado em voo é, portanto, produto de
uma malha de controle executada em tempo real a algumas centenas de hertz, o que faz
do quadrirrotor uma bancada didática particularmente completa: um único objeto de
estudo exige modelagem por dinâmica de corpo rígido, linearização em torno de um
ponto de operação, projeto de controlador, discretização e implementação embarcada
sob restrições de temporização.

Na prática, quase todos os quadrirrotores civis são governados por controladores
proporcional-integral-derivativo (PID) organizados em cascata, arranjo consagrado
pelo bom desempenho, pela robustez e pela facilidade de sintonia empírica. Os
principais projetos abertos de piloto automático --- entre eles o ArduPilot,
licenciado sob GPLv3 --- adotam essa estrutura: o comando é refinado sucessivamente
da posição para a velocidade, desta para a aceleração, para a atitude e, por fim,
para a taxa angular enviada à alocação de esforços dos motores. Trata-se de uma
arquitetura sintonizada por décadas de uso em campo, mas cujos ganhos são obtidos
majoritariamente por procedimentos de sintonia automática, e não a partir de um
modelo dinâmico explícito da aeronave. Surge daí a pergunta de pesquisa central
deste trabalho: \emph{um controlador projetado por métodos de controle moderno, a
partir de um modelo dinâmico obtido analiticamente e refinado por identificação, é
capaz de igualar ou superar o desempenho do PID nativo quando efetivamente embarcado
no firmware e submetido às mesmas condições de ensaio?}

Responder a essa pergunta exige mais do que uma simulação em ambiente de projeto.
Um controlador que se comporta bem em um simulador de propósito geral não é
equivalente a um controlador que se comporta bem dentro do laço de tempo real de um
piloto automático, onde a lei de controle convive com taxa de amostragem fixa,
atrasos de aquisição, filtragem de giroscópio, saturação de atuadores e comutação de
modos de voo. O ArduPilot oferece exatamente esse ponto de entrada: o
\emph{framework} \texttt{AC\_CustomControl} permite substituir a malha de atitude e
de taxa nativa por uma implementação própria em C++, invocada imediatamente antes da
biblioteca de motores, de modo a minimizar a latência introduzida, e comutável em
pleno voo por meio de um canal de rádio. Essa comutação é metodologicamente valiosa:
ela permite obter, em um mesmo ensaio e sob a mesma condição de vento e de
referência, o desempenho do controlador proposto e o da linha de base.

Há ainda um aspecto que distingue o desempenho aceitável em bancada do desempenho
aceitável em operação: a aeronave real não voa em ar parado. Vento médio, rajadas e
turbulência atmosférica constituem a principal fonte de perturbação externa sobre um
quadrirrotor de pequeno porte, e sua influência é agravada pela baixa inércia e pela
elevada relação entre área exposta e massa dessas aeronaves. A perturbação atua
simultaneamente como força sobre o centro de massa e como torque parasita nos eixos
de atitude, deslocando o ponto de operação em torno do qual o modelo foi linearizado.
Por esse motivo, avaliar um controlador apenas por sua resposta ao degrau em condição
nominal diz pouco sobre sua utilidade prática. Este trabalho incorpora, portanto, a
modelagem explícita de vento e de interferências externas tanto como \emph{cenário de
ensaio} --- por meio de perfis reprodutíveis de vento constante, rajada e turbulência
--- quanto como \emph{objeto de projeto}, mediante a inclusão de mecanismos de
estimação e compensação de distúrbio na estrutura do controlador. A capacidade de
rejeição de perturbação passa, assim, de efeito colateral a critério explícito de
comparação com a linha de base.

A validação ocorrerá em ambiente \emph{Software-In-The-Loop} (SITL), no qual o
firmware efetivamente compilado é executado contra um modelo de dinâmica de voo,
dispensando hardware e permitindo ensaios plenamente repetíveis, com distúrbios
programáveis e sem risco material ou de integridade física. Sob a ótica da
Engenharia de Controle e Automação, o trabalho percorre o ciclo completo de projeto
de um sistema de controle automático --- modelar a planta, identificá-la, projetar a
lei de controle, discretizá-la, embarcá-la e avaliá-la experimentalmente ---, o que
raramente se consegue reunir em um único trabalho de graduação com custo material
praticamente nulo.

A relevância do tema é reforçada por trabalhos consolidados na área. Bouabdallah
(2007) desenvolveu a modelagem, a identificação e o controle de um quadrirrotor,
comparando estratégias lineares e não-lineares em uma plataforma experimental
construída para esse fim. Sabatino (2015), por sua vez, derivou o modelo não-linear
de seis graus de liberdade e projetou controladores por realimentação de estados e
por técnicas não-lineares, avaliando-os exclusivamente em simulação numérica. Este
trabalho situa-se na lacuna entre essas duas abordagens: em vez de construir uma
aeronave própria --- o que desloca o esforço para a montagem mecânica e eletrônica
--- ou de encerrar a avaliação no ambiente de projeto --- o que deixa em aberto a
questão da implementação ---, a proposta consiste em \emph{embarcar o controlador
projetado em um piloto automático de uso real e de código aberto}, confrontando-o
diretamente com o controlador PID de produção sob condições de ensaio idênticas. Um
resultado adicional esperado é de natureza metodológica: documentar, em português,
um roteiro reprodutível que vai da equação de Newton-Euler até o código embarcado,
utilizando exclusivamente software livre.

% ----------------------------------------------------------------------
%  3. OBJETIVOS
% ----------------------------------------------------------------------
\section{Objetivos}

\subsection{Geral}

Projetar, implementar e validar em simulação um controlador de atitude alternativo
para um quadrirrotor, embarcado no firmware ArduPilot por meio do \emph{framework}
\texttt{AC\_CustomControl}, comparando quantitativamente seu desempenho ao do
controlador PID nativo em ambiente \emph{Software-In-The-Loop}, com uso exclusivo de
ferramentas de código aberto.

\subsection{Específicos}

\begin{itemize}[leftmargin=1.2cm,topsep=2pt]
    \item Obter o modelo dinâmico não-linear do quadrirrotor a partir das equações
          de Newton-Euler, incluindo a cinemática de atitude, o modelo dos
          atuadores e a matriz de alocação de esforços dos rotores, e levantar os
          parâmetros físicos correspondentes ao modelo simulado.
    \item Linearizar o modelo em torno da condição de voo pairado
          (\emph{hover}) por expansão em série de Taylor, obter a representação em
          espaço de estados contínua e discreta, verificar controlabilidade e
          observabilidade, e validar a linearização por comparação direta com o
          modelo não-linear.
    \item Caracterizar experimentalmente o controlador PID nativo do ArduPilot,
          estabelecendo-o como linha de base quantitativa de comparação.
    \item Identificar a dinâmica de atitude da aeronave simulada por meio do modo
          \emph{System Identification} do ArduPilot, estimando a resposta em
          frequência dos eixos e confrontando-a com o modelo analítico obtido.
    \item Projetar o controlador de atitude por técnicas de controle moderno ---
          adotando como abordagem principal o Regulador Linear Quadrático (LQR) e
          mantendo a alocação de polos como alternativa ---, verificar as margens de
          ganho e de fase resultantes e discretizá-lo para a taxa de execução do
          firmware.
    \item Modelar e implementar os distúrbios externos atuantes sobre a aeronave,
          contemplando vento médio constante, rajada determinística e turbulência
          estocástica segundo modelo espectral consagrado, bem como perturbações
          impulsivas de torque, de modo a compor cenários de ensaio reprodutíveis e
          parametrizáveis.
    \item Incorporar ao controlador um mecanismo de estimação e compensação de
          distúrbio --- por observador de estado estendido ou estrutura equivalente
          ---, avaliando o ganho de desempenho obtido em relação à realimentação de
          estados pura.
    \item Implementar o controlador em C++ no ArduPilot, com ganhos expostos como
          parâmetros ajustáveis em tempo de execução, registro dedicado dos sinais
          internos de referência, erro e esforço de controle, e tratamento adequado
          da reinicialização dos integradores na comutação entre controladores.
    \item Avaliar o desempenho em SITL por meio de ensaios de resposta ao degrau,
          rastreamento de referência, rejeição de distúrbio e variação paramétrica,
          comparando os resultados aos do PID nativo segundo métricas objetivas de
          desempenho transitório, de erro de rastreamento e de esforço de controle.
    \item Documentar todo o procedimento em formato reprodutível, com código e
          \emph{notebooks} versionados, de modo a permitir a continuidade do trabalho
          por outros alunos.
\end{itemize}

% ----------------------------------------------------------------------
%  4. METODOLOGIA
% ----------------------------------------------------------------------
\section{Metodologia}

O trabalho será conduzido em sete etapas encadeadas, descritas a seguir. A ordem é
deliberada: nenhuma linha de código embarcado será escrita antes de o controlador
estar validado em simulação numérica no ambiente de projeto, e nenhum ensaio
comparativo será realizado antes de a linha de base estar caracterizada.

\subsection{Modelagem}

O modelo dinâmico será derivado a partir das equações de Newton-Euler, adotando-se
o referencial inercial NED (\emph{North-East-Down}) e o referencial de corpo FRD
(\emph{Forward-Right-Down}), consistentes com a convenção empregada pelo ArduPilot.
A dinâmica de translação relaciona o empuxo total, orientado pela matriz de rotação,
à aceleração do centro de massa; a dinâmica de rotação relaciona os torques
aplicados à derivada da velocidade angular, incluindo o termo giroscópico
$\boldsymbol{\omega} \times \mathbf{I}\boldsymbol{\omega}$. A cinemática de atitude
relaciona as velocidades angulares medidas no corpo às derivadas dos ângulos de
Euler. O vetor de estado adotado reúne posição, velocidade, atitude e velocidade
angular, totalizando doze estados, e o vetor de entrada reúne o empuxo total e os
três torques. Cada rotor será modelado por leis quadráticas de empuxo e de torque de
reação, e a matriz de alocação mapeará os esforços desejados nas rotações
individuais dos rotores em função da geometria da aeronave. A derivação será
conduzida simbolicamente, de modo que a linearização da etapa seguinte seja obtida
diretamente por diferenciação. Efeitos de segunda ordem --- arrasto aerodinâmico do
corpo, \emph{blade flapping}, efeito solo e dinâmica de motor e ESC --- serão
inicialmente desprezados e justificados como fontes conhecidas de discrepância entre
modelo e planta.

\subsection{Linearização e análise}

A linearização será realizada em torno da condição de voo pairado, na qual os
ângulos de rolagem e arfagem são nulos, as velocidades linear e angular são nulas e
o empuxo equilibra o peso. As matrizes $\mathbf{A}$ e $\mathbf{B}$ serão obtidas
pelas matrizes Jacobianas de $\mathbf{f}(\mathbf{x},\mathbf{u})$ avaliadas nesse
ponto. Espera-se, como resultado estrutural, o desacoplamento aproximado do sistema
em subsistemas quase independentes de altitude, de atitude e de posição horizontal,
o que fundamenta tanto a arquitetura em cascata do firmware quanto a decisão de
projetar o controlador por malhas. Serão verificados o posto das matrizes de
controlabilidade e de observabilidade e os autovalores de $\mathbf{A}$, e o modelo
será discretizado por retentor de ordem zero na taxa de execução da malha interna. A
etapa se encerra com a validação da linearização: a resposta do modelo linear será
comparada à do modelo não-linear para perturbações de amplitude crescente,
evidenciando a faixa de validade da aproximação.

\subsection{Identificação do sistema}

O modo \emph{System Identification} do ArduPilot injeta sinais de excitação do tipo
\emph{chirp} nos eixos da aeronave e permite registrar a resposta correspondente. A
partir desses registros será estimada a resposta em frequência de cada eixo de
atitude e ajustada uma função de transferência equivalente por mínimos quadrados. A
comparação dessa estimativa com o modelo analítico da etapa anterior cumpre duas
funções: valida a modelagem e permite corrigir parâmetros mal estimados antes do
projeto do controlador, reduzindo o risco de projetar sobre um modelo incorreto.

\subsection{Modelagem de vento e interferências externas}

Os distúrbios externos serão tratados como parte integrante do modelo, e não como
ruído residual. Adota-se a decomposição usual do campo de vento em três componentes.
A primeira é o \emph{vento médio}, representado por um vetor de velocidade constante
no referencial inercial, definido por intensidade e direção; sua ação é
predominantemente de regime permanente e exige do controlador ação integral ou
compensação explícita. A segunda é a \emph{rajada}, modelada de forma determinística
pelo perfil clássico de meio ciclo de cosseno, parametrizado por amplitude e
comprimento característico, adequado à avaliação da resposta transitória a uma
perturbação de curta duração e amplitude conhecida. A terceira é a \emph{turbulência},
modelada como processo estocástico obtido pela filtragem de ruído branco por funções
de transferência de conformação espectral, segundo o modelo de Dryden, cujos
parâmetros de intensidade e de escala de comprimento são tabelados em função da
altitude e da severidade da condição atmosférica.

A conversão dessas componentes em esforços sobre a aeronave se dá pela velocidade
relativa entre o ar e o corpo. A força de arrasto resultante atua sobre o centro de
massa e, quando o centro de pressão não coincide com este, produz um torque parasita
nos eixos de atitude --- efeito que, no modelo nominal da Seção~4.1, havia sido
deliberadamente desprezado e que aqui é reintroduzido de maneira controlada, na forma
de um termo aditivo de perturbação na equação da dinâmica de rotação. Além do vento,
serão consideradas interferências externas de outras naturezas, tratadas sob a mesma
formulação de perturbação aditiva: variação abrupta de massa ou de posição do centro
de gravidade, perda parcial de eficiência de um rotor, e perturbações impulsivas de
torque, empregadas para avaliar a recuperação do controlador a partir de condições
iniciais deslocadas.

Do ponto de vista de implementação, os cenários serão reproduzidos por dois caminhos
complementares. No ambiente de projeto em Python, os modelos de vento serão
incorporados diretamente à simulação do modelo não-linear, permitindo o projeto e a
avaliação preliminar da rejeição de perturbação. Em SITL, serão utilizados os
parâmetros nativos de simulação de vento do ArduPilot, que permitem configurar
intensidade, direção, turbulência e rajada sem alteração de código, garantindo que os
ensaios comparativos ocorram sob perfis idênticos para ambos os controladores. A
consistência entre os dois caminhos --- modelo em Python e simulação em SITL --- será
verificada previamente aos ensaios definitivos.

\subsection{Projeto do controlador}

O controlador será projetado sobre o subsistema de atitude do modelo linearizado.
Na abordagem principal, o ganho de realimentação de estados será obtido pela solução
da equação algébrica de Riccati associada ao problema do Regulador Linear
Quadrático, com as matrizes de ponderação $\mathbf{Q}$ e $\mathbf{R}$ ajustadas
iterativamente a partir da resposta em malha fechada desejada. Serão avaliados o
desempenho transitório em resposta ao degrau, o rastreamento de referências variantes
no tempo e as margens de ganho e de fase resultantes. A alocação de polos será
mantida como abordagem alternativa, de valor didático comparativo. Caso a
disponibilidade de estados medidos se mostre insuficiente ou o ruído de medição se
mostre relevante, prevê-se a extensão para uma estrutura LQG, com estimador de
estados.

O projeto contemplará explicitamente a rejeição das perturbações descritas na
subseção anterior. A realimentação de estados pura, por si só, não anula erro em
regime permanente diante de um distúrbio constante, situação exatamente
correspondente ao vento médio. Serão avaliadas, por isso, duas estratégias
complementares. A primeira é a inclusão de ação integral por aumento do vetor de
estados com o integral do erro de atitude, o que preserva a formulação quadrática de
projeto e recai em um problema de mesma natureza, acrescido do tratamento de
saturação do termo integral. A segunda é a estimação ativa do distúrbio por meio de
um observador de estado estendido, no qual o torque perturbador é tratado como estado
adicional a ser estimado e subtraído do sinal de controle, caracterizando uma
compensação por antecipação. As duas estratégias serão comparadas entre si e com a
realimentação pura, de modo a quantificar o ganho efetivo atribuível à compensação de
distúrbio. Antes da implementação, a lei de controle resultante será discretizada na
taxa de execução da malha interna e reavaliada no domínio discreto, verificando-se em
particular a sensibilidade do observador ao ruído de medição do giroscópio.

\subsection{Implementação embarcada}

A implementação partirá do \emph{backend} de exemplo fornecido pelo próprio
\emph{framework} \texttt{AC\_CustomControl}, o que restringe o esforço de programação
à codificação da lei de controle propriamente dita e ao tratamento de estado interno.
A função de atualização receberá a atitude e a taxa angular de referência já
processadas pela cadeia nativa, bem como a medida mais recente do giroscópio, e
devolverá os esforços de rolagem, arfagem e guinada à alocação de motores. Três
cuidados serão observados: a limitação do termo integral contra o efeito de
\emph{windup}; a reinicialização correta dos estados internos na comutação entre
controladores, de modo a evitar transitórios espúrios; e a exposição dos ganhos como
parâmetros do firmware, permitindo sintonia em tempo de execução pela estação de
controle em solo, sem recompilação. Serão adicionadas mensagens de registro
dedicadas para os sinais internos de cada eixo, planejadas antecipadamente, uma vez
que são elas que fundamentarão a análise de resultados.

\subsection{Ensaios e métricas}

A bateria de ensaios em SITL compreenderá quatro grupos. O primeiro, em condição
nominal de ar parado, avaliará a resposta ao degrau em rolagem e arfagem e o
rastreamento de trajetórias de atitude e de posição, caracterizando o desempenho
transitório básico. O segundo, de rejeição de perturbação, submeterá a aeronave aos
cenários de vento definidos na Seção~4.5, organizados em níveis crescentes de
severidade: vento médio constante em diferentes intensidades e direções relativas ao
eixo de comando; rajada isolada de amplitude e duração conhecidas; e turbulência
contínua em condições leve, moderada e severa. O terceiro grupo tratará das demais
interferências externas, com perturbação impulsiva de torque, variação abrupta de
massa e perda parcial de eficiência de rotor. O quarto avaliará robustez a
incertezas de modelagem por variação paramétrica de massa e de momentos de inércia
do modelo simulado.

Cada ensaio será executado com o controlador proposto, em suas variantes com e sem
compensação de distúrbio, e com o PID nativo, sob perfis de perturbação idênticos e
sementes idênticas nos casos estocásticos, o que assegura comparabilidade direta. Os
registros de voo serão processados computacionalmente e as seguintes métricas serão
tabeladas: tempo de subida, sobressinal, tempo de acomodação e erro em regime
permanente, para os ensaios transitórios; desvio máximo de atitude, tempo de
recuperação e raiz do erro quadrático médio, para os ensaios sob perturbação; e
esforço de controle acumulado, em todos os casos, como indicador do custo energético
associado ao desempenho obtido. A discussão dos resultados confrontará as diferenças
observadas com as previsões da teoria, com as limitações conhecidas do modelo e com o
compromisso entre agressividade de rejeição e esforço de atuação.

% ----------------------------------------------------------------------
%  5. INFORMAÇÕES COMPLEMENTARES
% ----------------------------------------------------------------------
\section{Informações Complementares}

\subsection{Recursos e ferramentas}

O projeto será desenvolvido integralmente com ferramentas de código aberto,
organizadas em duas camadas complementares: uma camada de projeto e análise, em
Python, e uma camada de tempo real, em C++, dentro do próprio firmware. A
Tabela~\ref{tab:ferramentas} resume os recursos previstos. Não há custo de aquisição
de material, uma vez que toda a validação ocorre em simulação.

\begin{table}[htbp]
\centering
\small
\caption{Ferramentas previstas e respectivos papéis no trabalho.}
\label{tab:ferramentas}
\begin{tabular}{|L{4.2cm}|L{7.3cm}|L{2.6cm}|}
\hline
\textbf{Ferramenta} & \textbf{Papel no trabalho} & \textbf{Camada}\\
\hline
Ubuntu (nativo ou WSL2) e Git & Ambiente de desenvolvimento e controle de versão & Base\\
\hline
ArduPilot (GPLv3) & Firmware de piloto automático; hospeda o controlador proposto & C++\\
\hline
SITL e \texttt{sim\_vehicle.py} & Execução do firmware compilado sem hardware & Simulação\\
\hline
MAVProxy e QGroundControl & Estações de controle em solo; comando e ajuste de parâmetros & Telemetria\\
\hline
\texttt{sympy} & Derivação simbólica do modelo e linearização por Jacobiana & Python\\
\hline
\texttt{python-control} & Projeto do controlador e análise de estabilidade & Python\\
\hline
\texttt{numpy} e \texttt{scipy} & Álgebra linear, integração numérica e identificação & Python\\
\hline
Modelos de vento (Dryden e rajada) e parâmetros de simulação de vento do SITL & Geração dos cenários de perturbação atmosférica reprodutíveis & Simulação\\
\hline
\texttt{pymavlink}, \texttt{pandas} e \texttt{matplotlib} & Leitura dos registros de voo, métricas e gráficos & Python\\
\hline
Gazebo com \texttt{ardupilot\_gazebo} & Física e visualização tridimensional (opcional) & Simulação\\
\hline
\end{tabular}
\end{table}

\subsection{Delimitação do escopo}

O trabalho restringe-se à malha de atitude --- rolagem, arfagem e guinada --- com
possível extensão à malha de altitude caso o cronograma permita, e à validação em
ambiente de simulação. Quanto às perturbações, o escopo abrange os modelos de vento
médio, rajada determinística e turbulência estocástica aplicados de forma uniforme ao
volume ocupado pela aeronave, além de perturbações aditivas de torque e de variação
paramétrica; ficam fora do escopo a modelagem de campos de vento espacialmente
distribuídos, esteiras de obstáculos e efeitos aerodinâmicos de proximidade ao solo
ou a paredes, que exigiriam simulação fluidodinâmica. Ficam explicitamente fora do escopo o voo em hardware real, o
controle de trajetória por métodos preditivos, a operação cooperativa de múltiplas
aeronaves e o uso de visão computacional, todos registrados como possibilidades de
trabalhos futuros. Caso haja folga significativa de cronograma, prevê-se, como
extensão opcional, a execução em modo \emph{Simulation on Hardware}, que mantém a
simulação da planta mas executa o firmware na controladora de voo real.

\subsection{Riscos e mitigações}

Os principais riscos identificados são de natureza técnica e de escopo. O primeiro é
a configuração do ambiente de compilação do firmware, etapa reconhecidamente
sensível e capaz de consumir tempo desproporcional; como mitigação, ela será
concluída integralmente na fase inicial, antes de qualquer atividade de modelagem, e
tomada como marco de entrada para as demais fases. O segundo é a divergência entre o
modelo analítico e a dinâmica efetivamente simulada, mitigada pela etapa de
identificação de sistemas, que permite corrigir os parâmetros antes do projeto do
controlador. O terceiro é a ocorrência de transitórios indesejados na comutação entre
controladores, mitigada pelo tratamento explícito da reinicialização dos estados
internos e por ensaios específicos de comutação. O quarto é a possibilidade de o
controlador projetado não superar o PID nativo em todas as métricas; entende-se, no
entanto, que esse não é um risco ao trabalho, mas um resultado legítimo, desde que a
comparação seja metodologicamente correta e as causas sejam analisadas à luz da
teoria. O quinto é a expansão indevida do escopo, mitigada pela delimitação
apresentada na subseção anterior. O sexto diz respeito à representatividade dos
modelos de distúrbio: parâmetros de turbulência excessivamente severos podem levar à
saturação dos atuadores e tornar a comparação entre controladores pouco informativa,
enquanto parâmetros brandos demais podem mascarar diferenças relevantes; como
mitigação, os níveis de severidade serão definidos previamente a partir de valores
tabelados na literatura e verificados quanto à ocorrência de saturação antes da
execução da bateria definitiva de ensaios.

% ----------------------------------------------------------------------
%  6. CRONOGRAMA PREVISTO
% ----------------------------------------------------------------------
\section{Cronograma Previsto}

O cronograma da Tabela~\ref{tab:cronograma} distribui as atividades ao longo de
aproximadamente vinte e quatro semanas. As fases se sobrepõem deliberadamente, de
modo que o estudo teórico de cada etapa ocorra durante a execução da etapa anterior.

\begin{table}[htbp]
\centering
\small
\caption{Cronograma previsto de execução.}
\label{tab:cronograma}
\begin{tabular}{|c|L{6.0cm}|c|L{4.4cm}|}
\hline
\textbf{Fase} & \textbf{Atividade principal} & \textbf{Semanas} & \textbf{Entregável}\\
\hline
0 & Preparação do ambiente, compilação do firmware e primeiro voo simulado & 1--2 & Ambiente funcional e leitura de registros\\
\hline
1 & Modelagem dinâmica não-linear e levantamento de parâmetros & 3--6 & Modelo simbólico documentado\\
\hline
2 & Linearização, análise estrutural e discretização & 5--8 & Modelo em espaço de estados validado\\
\hline
3 & Caracterização do PID nativo e projeto do controlador & 8--12 & Controlador projetado e simulado\\
\hline
4 & Identificação do sistema, refino do modelo e implementação dos modelos de vento e de perturbação & 10--15 & Modelo identificado e cenários de distúrbio validados\\
\hline
5 & Projeto da compensação de distúrbio e implementação em C++ no firmware & 13--18 & Controlador operante em SITL\\
\hline
6 & Ensaios comparativos sob condição nominal e sob perturbação, métricas e análise & 17--22 & Resultados e gráficos comparativos\\
\hline
7 & Redação e fechamento da monografia & 21--24 & Monografia finalizada\\
\hline
\end{tabular}
\end{table}

% ----------------------------------------------------------------------
%  7. APLICAÇÃO DOS CONHECIMENTOS
% ----------------------------------------------------------------------
\section{Aplicação dos Conhecimentos}

O desenvolvimento deste projeto integra conhecimentos fundamentais adquiridos ao
longo do curso de Engenharia de Controle e Automação, incluindo disciplinas como
Mecânica Aplicada e Dinâmica de Corpos Rígidos, Cálculo e Álgebra Linear, Modelagem
e Simulação de Sistemas Dinâmicos, Sistemas de Controle Clássico e Moderno, Controle
Digital, Identificação de Sistemas, Sistemas Microprocessados e Embarcados, Sistemas
em Tempo Real e Programação. Esses conhecimentos são essenciais para a derivação do
modelo dinâmico da aeronave, para o projeto e a discretização da lei de controle,
para a sua implementação dentro de um laço de execução de tempo real e para a
interpretação crítica dos resultados experimentais obtidos em simulação.

A articulação entre essas áreas é o que confere unidade ao trabalho. A dinâmica de
corpo rígido fornece as equações de movimento; a álgebra linear e o cálculo
sustentam a linearização por matrizes Jacobianas e a análise por autovalores; o
controle moderno fornece o critério de projeto e a formulação por realimentação de
estados; o controle digital fornece a discretização e a análise dos efeitos de
amostragem; a identificação de sistemas fornece o instrumento de validação do modelo;
e os conhecimentos de sistemas embarcados e de programação sustentam a implementação
sob restrição de temporização. O resultado é um exercício de engenharia de controle
em sua forma integral, e não apenas o estudo isolado de uma técnica.

Para embasar o desenvolvimento serão utilizadas as referências clássicas de
modelagem e controle de quadrirrotores de Bouabdallah~\cite{bouabdallah2007},
Sabatino~\cite{sabatino2015} e Beard~\cite{beard2008}, que fornecem a formulação por
Newton-Euler, a linearização em torno do voo pairado e o projeto por realimentação
de estados. A fundamentação em teoria de controle apoiar-se-á nas obras de Åström e
Murray~\cite{astrom2008}, base conceitual da biblioteca \texttt{python-control}, e de
Ogata~\cite{ogata2010} para os aspectos de espaço de estados e de projeto por
critérios quadráticos, complementadas por Franklin, Powell e
Workman~\cite{franklin1998} no que se refere à discretização e à implementação
digital, e por Ljung~\cite{ljung1999} no que se refere à identificação de sistemas. A
modelagem dos distúrbios atmosféricos apoiar-se-á em Beard e McLain~\cite{beard2012},
que formulam o modelo de vento aplicado a veículos aéreos não tripulados de pequeno
porte, e na especificação militar norte-americana~\cite{mil8785c}, origem dos
espectros de turbulência e dos respectivos parâmetros tabelados; a estratégia de
compensação ativa de perturbação apoiar-se-á na formulação de controle por rejeição
ativa de distúrbio de Han~\cite{han2009}.
A especificação da arquitetura de controle em cascata, da interface de controlador
customizado e do modo de identificação de sistemas será extraída da documentação
oficial de desenvolvimento do ArduPilot~\cite{ardupilotdev}, que servirá de
referência tanto para a implementação embarcada quanto para a metodologia de ensaio
e análise dos resultados.

% ----------------------------------------------------------------------
%  8. REFERÊNCIAS
% ----------------------------------------------------------------------
\section*{Referências}
\addcontentsline{toc}{section}{Referências}

\begin{thebibliography}{9}
\setlength{\itemsep}{0.4em}

\bibitem{bouabdallah2007}
BOUABDALLAH, Samir. \textbf{Design and control of quadrotors with application to
autonomous flying}. 2007. Tese (Doutorado) -- École Polytechnique Fédérale de
Lausanne (EPFL), Lausanne, 2007.

\bibitem{sabatino2015}
SABATINO, Francesco. \textbf{Quadrotor control: modeling, nonlinear control design,
and simulation}. 2015. Dissertação (Mestrado em Engenharia Elétrica) -- KTH Royal
Institute of Technology, Estocolmo, 2015.

\bibitem{beard2008}
BEARD, Randal W. \textbf{Quadrotor dynamics and control}. Provo: Brigham Young
University, 2008. (Relatório técnico).

\bibitem{astrom2008}
ÅSTRÖM, Karl J.; MURRAY, Richard M. \textbf{Feedback systems: an introduction for
scientists and engineers}. Princeton: Princeton University Press, 2008.

\bibitem{ogata2010}
OGATA, Katsuhiko. \textbf{Engenharia de controle moderno}. 5. ed. São Paulo:
Pearson Prentice Hall, 2010.

\bibitem{franklin1998}
FRANKLIN, Gene F.; POWELL, J. David; WORKMAN, Michael L. \textbf{Digital control of
dynamic systems}. 3rd ed. Menlo Park: Addison-Wesley, 1998.

\bibitem{ljung1999}
LJUNG, Lennart. \textbf{System identification: theory for the user}. 2nd ed. Upper
Saddle River: Prentice Hall, 1999.

\bibitem{beard2012}
BEARD, Randal W.; McLAIN, Timothy W. \textbf{Small unmanned aircraft: theory and
practice}. Princeton: Princeton University Press, 2012.

\bibitem{mil8785c}
UNITED STATES DEPARTMENT OF DEFENSE. \textbf{MIL-F-8785C: military specification --
flying qualities of piloted airplanes}. Washington, D.C.: Department of Defense,
1980.

\bibitem{han2009}
HAN, Jingqing. From PID to active disturbance rejection control. \textbf{IEEE
Transactions on Industrial Electronics}, v. 56, n. 3, p. 900--906, 2009.

\bibitem{ardupilotdev}
ARDUPILOT DEVELOPMENT TEAM. \textbf{ArduPilot developer documentation}. Disponível
em: \url{https://ardupilot.org/dev/}. Acesso em: [data de acesso].

\end{thebibliography}

\end{document}